\subsection{Six exact genus-two descents and the rank-fourteen Jacobian}

\begin{statement}{A.186}{Proved}
{All six genus-two ranks, their two-primary Tate--Shafarevich groups,
and the exact rank of the pointed Fermat-eight Jacobian}
For every pair in the set \(\mathcal G\) of (A.182.3), put
\(J_{A,B}=\operatorname{Jac}(G_{A,B})\).  Then
\[
\begin{aligned}
 \operatorname{rank}J_{A,B}(\Q)&=1,
 &J_{A,B}(\Q)_{\rm tors}&\simeq\Z/2,\\
 \operatorname{Sel}^{2}(J_{A,B}/\Q)&\simeq(\Z/2)^2,
 &\Sha(J_{A,B}/\Q)[2^\infty]&=0.
\end{aligned}
\tag{A.186.1}
\]
The exact two-Selmer images are the explicit two-dimensional
kernels in (A.186.16) below, under the literal odd-degree Kummer
maps.  For every \(n\geq1\), the finite Kummer map identifies
\[
 J_{A,B}(\Q)/2^nJ_{A,B}(\Q)
       \xrightarrow{\ \sim\ }
       \operatorname{Sel}^{2^n}(J_{A,B}/\Q).
\tag{A.186.2}
\]
For \(c=A^3B\), put
\(J_c=\operatorname{Jac}(C_c)\),
\(C_c:y^2=x^5+cx\), using the actual coordinate
isomorphism (A.186.4).  The literal point classes
\(D_c\) and the rational
branch classes \(T_c\) displayed in (A.186.18)
give the actual topological isomorphisms
\(\Z_2\oplus\Z/2\to\widehat{J_c(\Q)}_2\),
\((a,e)\mapsto aD_c+eT_c\).
Together with A.185 and the actual decomposition in A.182, this
gives the exact Mordell--Weil rank
\[
                 \operatorname{rank}J(C_8)(\Q)=14<21=g(C_8).
\tag{A.186.3}
\]
The proof supplies the full global squareclass bounds and seventeen
actual local point tests.  It does not deduce a classification of
sections or a description of a local Chabauty zero set from (A.186.3).
\end{statement}

\begin{proof}
\emph{The curve isomorphisms and the coefficient module.}
The actual change of coordinates
\[
 x=AX,\qquad y=A^2Y
\tag{A.186.4}
\]
identifies \(G_{A,B}:Y^2=AX^5+BX\) with the proper smooth model
\(C_c:y^2=x^5+cx\), where \(c=A^3B\).  Its inverse is
\((X,Y)=(x/A,y/A^2)\), and both maps fix the unique point
\(\infty\).  In the order of \(\mathcal G\) used in A.182,
the six coefficients are
\[
                (c_1,\ldots,c_6)
                =(-54,648,-432,13824,27,-4608).
\tag{A.186.5}
\]
Write \(F_c(T)=T^4+c\), \(K_c=\Q[\theta]/(F_c)\), and
\(L_c=\Q\times K_c\).  Each \(F_c\) is irreducible.
For \(c=-54,-432,13824,27\), its four roots have
three-adic valuation \(3/4\); for \(648\) and \(-4608\)
their two-adic valuations are \(3/4\) and \(9/4\).
The constant term of a monic local factor has integral valuation,
forcing the factor degree to be divisible by four.  This proves
irreducibility over the indicated local field and over \(\Q\).
Consequently \(xF_c(x)\) has exactly two rational irreducible
factors.

Here are the coefficient and pairing identifications, including
the rational factor.  Label the five finite branch points by
\(P_i=(\alpha_i,0)\), and put \(D_i=[P_i-\infty]\).
These generate \(J_c[2]\), with the single relation
\(\sum_iD_i=0\).  The usual divisor functions
\(x-\alpha_i\), whose divisors are \(2D_i\), give
\(e_2(D_i,D_j)=-1\) for \(i\ne j\) and
\(e_2(D_i,D_i)=1\).  Thus the actual pairing map
\[
 J_c[2]\longrightarrow
 M_c:=\ker\bigl(\operatorname{Res}_{L_c/\Q}\mu_2
                         \xrightarrow{N}\mu_2\bigr),
 \qquad D\longmapsto(e_2(D,D_i))_{i=1}^5
\tag{A.186.6}
\]
is a Galois-equivariant isomorphism.  The image of \(D_i\)
has four minus signs and a plus sign at position \(i\).
The coordinate dot pairing on \(M_c\) gives the asserted
Weil pairing on these generators: intersections have size
three for distinct indices and four for equal indices.
Equivalently, in additive coordinates the actual map from
\(\F_2^5/\langle(1,1,1,1,1)\rangle\) to
\(\ker(\sum:\F_2^5\to\F_2)\) is
\(u\mapsto u+(\sum u)(1,1,1,1,1)\).

Since five is odd, the diagonal \(\mu_2\) splits the norm
map on the full five-coordinate module.  Kummer theory and
Shapiro therefore identify its cohomology exactly as
\[
\begin{aligned}
 H^1(k,J_c[2])
 &\simeq\ker\bigl(L_{c,k}^{\times}/L_{c,k}^{\times2}
                    \xrightarrow{N}k^\times/k^{\times2}\bigr),\\
 K_{c,k}^{\times}/K_{c,k}^{\times2}
 &\xrightarrow{\ \sim\ }H^1(k,J_c[2]),
 &\eta&\longmapsto(N_{K_{c,k}/k}\eta,\eta).
\end{aligned}
\tag{A.186.7}
\]
Here \(K_{c,k}=K_c\otimes_\Q k\) can be a product of
fields.  The first coordinate of any norm-kernel class is
forced by its second coordinate; inverses and the original
element have the same squareclass.  In particular there is
no further quotient by the diagonal \(k^\times\).
Changing any representative by a square changes neither map
in (A.186.7).

The actual odd-degree descent map sends a point divisor
\([P-\infty]\), for \(P=(q,y)\) off the branch locus,
to
\[
                (q,q-\theta)\in L_{c,k}^\times/L_{c,k}^{\times2}.
\tag{A.186.8}
\]
Its norm is \(q(q^4+c)=y^2\).  We verify directly that
this is the finite Kummer map under (A.186.6).
Represent a class by a degree-zero divisor \(D\) over
\(k\), moved off the branch points and infinity; the
rational point \(\infty\) permits such a divisor
representative.  Over \(\overline k\), choose a divisor
\(D'\) representing a half of its class, with the same
disjointness, and a function \(g\) with
\(\operatorname{div}(g)=2D'-D\).  Write
\(f_i=x-\alpha_i\), and use \(D_i=P_i-\infty\)
as an actual divisor in these evaluations.  Weil
reciprocity gives
\[
 f_i(D')^2/f_i(D)=g(D_i)^2,\qquad
 \beta_i:=f_i(D')/g(D_i),\qquad \beta_i^2=f_i(D).
\]
For \(\sigma\in G_k\), put
\(E_\sigma=\sigma D'-D'\) and
\(h_\sigma=\sigma g/g\), so that
\(\operatorname{div}(h_\sigma)=2E_\sigma\).
With the natural permutation action on the five
coordinates, the square-root cocycle is
\[
 \frac{(\sigma\beta)_i}{\beta_i}
   =\frac{f_i(E_\sigma)}{h_\sigma(D_i)}
   =e_2([E_\sigma],D_i).
\]
The two possible reciprocal conventions for the Weil
pairing agree in \(\mu_2\).  The class
\([E_\sigma]\) is precisely the finite Kummer
cocycle.  Thus these literal function evaluations give
the coefficient identification (A.186.6), and their
kernel is \(2J_c(k)\).

For evaluation on \(P-\infty\), the local parameter
\(t=x^2/y\) at infinity gives
\((x-\alpha_i)t^2\to1\); its normalized leading
coefficient contributes one.  This proves (A.186.8).
At \(P_0=(0,0)\), the local parameter \(y\) gives
\(x/y^2\to1/c\).  Since \(1/c\) and \(c\)
have the same squareclass, the rational branch divisor
\(T_0=[P_0-\infty]\) has coordinate
\((c,-\theta)\).  These computations also give the
point and branch formulas of
\cite[Lemma 4.1 and Lemma 4.3(1)--(2), pp.250--251]{LegStollDescent2001}
with exactly the same coefficient map.

\emph{The exact local linear conditions.}
The coordinate dot pairing in (A.186.6) makes the local
Tate pairing the product of quadratic Hilbert symbols on
the components of \(L_{c,\Q_p}\).  Pairing the class
\((N\eta,\eta)\) with (A.186.8) gives
\[
\begin{aligned}
 (N\eta,q)_p\prod_{\mathfrak p\mid p}
                    (\eta,q-\theta)_{\mathfrak p}
 &=\prod_{\mathfrak p\mid p}
                    (\eta,q(q-\theta))_{\mathfrak p}.
\end{aligned}
\tag{A.186.9}
\]
This is the norm/corestriction projection formula for the
first factor.  At \(T_0\) the corresponding second argument
is \(-c\theta\).

The image of \(J_c(\Q_p)/2J_c(\Q_p)\) is isotropic
for this pairing, including at \(p=2\).  To specify the
polarization precisely, the image of
\(C_c\to J_c,\ P\mapsto[P-\infty]\), is the rational
symmetric theta divisor: inversion on it is induced by the
hyperelliptic involution.  Its principal polarization
identifies \(J_c\) with its dual, and \(\mathcal O(2\Theta)\)
induces twice that polarization, hence the isogeny \([2]\)
in (A.186.7).  The Kummer isotropy and its Weil-cup pairing
are therefore precisely those in
\cite[Section 4.1, Corollary 4.7, and Propositions 4.9 and 4.11,
pp.16--18]{LegPoonenRains2012}.  It follows that every
two-Selmer class satisfies (A.186.9) with value one for
every actual local point used below.  Only these necessary
conditions are used; the selected points need not generate
the entire local Mordell--Weil group.

\emph{The complete global candidate group.}
The polynomial discriminant is literally
\[
                       \operatorname{disc}(x^5+cx)=256c^5.
\tag{A.186.10}
\]
Thus a two-Selmer class has even valuations at every prime
of \(K_c\) outside two and three, by
\cite[Proposition 4.6 and Corollary 4.7, pp.252--253]{LegStollDescent2001}.
Let \(S_c\) be all primes of \(K_c\) above two and three.
The full number-field computations used here give
\[
\begin{array}{r|r|c|c|c|r|r}
c&\operatorname{disc}K_c&(r_1,r_2)&\operatorname{Cl}(K_c)
 &\operatorname{Cl}(\mathcal O_{K_c}[1/6])&|S_c|&d_c\\ \hline
-54&-55296&(2,1)&0&0&2&5\\
648&2048&(0,2)&0&0&4&6\\
-432&-6912&(2,1)&0&0&2&5\\
13824&55296&(0,2)&\Z/2&0&2&4\\
27&432&(0,2)&0&0&2&4\\
-4608&-18432&(2,1)&0&0&2&5.
\end{array}
\tag{A.186.11}
\]
Here \(d_c=r_1+r_2-1+|S_c|+1\).  In the only nontrivial
ordinary class group, both the prime above two and the
prime above three represent its nonzero class.
The localized class group being zero, every squareclass
with even valuations outside \(S_c\) has an \(S_c\)-unit
representative.  Indeed, the half-divisor outside \(S_c\)
becomes principal in the localized ideal group; dividing
by the square of its principal generator removes it.
The global candidate group is consequently exactly
\(\mathcal O_{K_c,S_c}^{\times}/
          \mathcal O_{K_c,S_c}^{\times2}\), of dimension \(d_c\).

We give all literal generators used in the computations.
In each row \((1,u,v,w)\) is the displayed integral basis
of \(\mathcal O_{K_c}\):
\[
\begin{array}{r|c|c|c}
c&u&v&w\\ \hline
-54&\theta^3/9&\theta^2/3&\theta\\
648&\theta^3/108&\theta^2/18&\theta/3\\
-432&\theta^3/72&-\theta^2/12&-\theta/2\\
13824&\theta^3/576&\theta^2/48&\theta/4\\
27&(3-\theta^2)/6&(9\theta-\theta^3)/18
                       &(\theta^3+9\theta)/18\\
-4608&\theta^2/48&-\theta/4&\theta^3/192.
\end{array}
\tag{A.186.12}
\]
The ordered squareclass bases \(G_c=(g_1,\ldots,g_{d_c})\)
are the following elements, with the row-specific meaning
of \(u,v,w\):
\[
\begin{aligned}
G_{-54}={}&(u-v+w-2,\ 3-u+v-w,\ 5+2v,\
                         5-2u+2v-2w,\ -1),\\
G_{648}={}&(u,\ u-1,\ u+1,\ 1-v,\ -1-v-w,\ -1),\\
G_{-432}={}&(u+1,\ u,\ 2-v,\ 2+u-v-w,\ -1),\\
G_{13824}={}&(10-4u-v+2w,\ -6+5u+4v+w,\
                         -95-46u+20v+12w,\ -1),\\
G_{27}={}&(1-u+v,\ v,\ 1-u-w,\ u),\\
G_{-4608}={}&(2+2u-v+w,\ -3-3u+2v-w,\ -1-u,\
                         17+12u-8v+6w,\ -1).
\end{aligned}
\tag{A.186.13}
\]
The retained certificate constructs each field with its
displayed defining polynomial and invokes
\texttt{bnfcertify} with its default flag zero, which
certifies the full class group and full unit group.
It computes the localized class group, obtains a full
fundamental \(S_c\)-unit system, and expresses every
literal element of (A.186.13) in that system using
\texttt{bnfisunit}.  The resulting square matrices have
full rank \(d_c\) modulo two.  Thus the listed elements
span every candidate squareclass, not merely a subgroup
of finite index of unspecified parity.  The meanings and
unconditional certification guarantee of these specific
routines are given in the official primary documentation
\cite[entries \texttt{bnfcertify}, \texttt{bnfsunit},
\texttt{bnfunits}, and \texttt{bnfisunit}]{LegPariNumberFields}.
All integral bases, prime ideals, class data, and full
coordinate matrices are retained in the executed output.

\emph{The seventeen actual local points and all rows.}
For each listed pair \((p,q)\), let
\(\nu=v_p(q^5+cq)\).  The following table gives \(\nu\)
in the same order as the point coordinates:
\[
\begin{array}{r|l|l}
c&(p,q)&\nu\\ \hline
-54&(2,-96),(2,-93),(2,-86)&6,0,2\\
648&(2,-95),(2,-90),(2,-56),(3,-96)&0,4,6,6\\
-432&(2,-95),(2,-90),(2,-76)&0,6,6\\
13824&(2,-100),(2,-40)&10,12\\
27&(2,-97),(2,-93)&2,2\\
-4608&(2,-100),(2,-72),(2,-32)&10,12,14.
\end{array}
\tag{A.186.14}
\]
Every \(\nu\) is even.  In every dyadic row
\((q^5+cq)/2^\nu\equiv1\pmod8\); in the one
three-adic row the corresponding unit is one modulo three.
These are exact rational congruences.  Hensel lifting gives
an actual point
\(P_{p,q}=(q,p^{\nu/2}\sqrt{(q^5+cq)/p^\nu})\)
of \(C_c(\Q_p)\).  One can choose the dyadic square root
congruent to one modulo four, and the three-adic one
congruent to one modulo three.  None of these points is
a branch point.

Use row order (A.186.14), column order (A.186.13), and
encode a Hilbert symbol of \(-1\) by one in \(\F_2\).
The complete matrices from (A.186.9) are
\[
\begin{aligned}
M_{-54}&=\begin{pmatrix}1&0&1&1&0\\0&0&1&1&1\\0&1&1&0&0\end{pmatrix},
&M_{648}&=\begin{pmatrix}0&1&0&1&0&0\\1&1&0&0&0&1\\1&1&0&0&0&0\\0&0&0&0&1&1\end{pmatrix},\\
M_{-432}&=\begin{pmatrix}1&0&0&1&0\\0&0&1&1&1\\0&1&1&0&0\end{pmatrix},
&M_{13824}&=\begin{pmatrix}0&1&1&1\\0&0&0&1\end{pmatrix},\\
M_{27}&=\begin{pmatrix}0&1&0&1\\1&1&1&1\end{pmatrix},
&M_{-4608}&=\begin{pmatrix}0&1&0&1&1\\1&1&1&0&1\\0&0&0&0&1\end{pmatrix}.
\end{aligned}
\tag{A.186.15}
\]
The executed certificate prints each individual symbol at
each prime above \(p\), multiplies them exactly as in
(A.186.9), and asserts equality with every displayed entry.
The local Hilbert routine is the quadratic conic/norm
symbol specified in the primary documentation
\cite[entry \texttt{nfhilbert}]{LegPariNumberFields}.
In particular the last row for \(c=648\) includes all
three primes above three, of residue degrees one, one,
and two; none is omitted.

The matrix ranks are \(3,4,3,2,2,3\), respectively.
All kernels therefore have dimension two.  To give the
exact surviving classes, the following two binary rows
are bases of those kernels; a row \(e\) means
\(\eta=\prod_jg_j^{e_j}\), followed by the actual map
\(\eta\mapsto(N\eta,\eta)\) of (A.186.7):
\[
\begin{array}{r|c|c}
c&e^{(1)}&e^{(2)}\\ \hline
-54&(0,1,1,1,0)&(1,1,1,0,1)\\
648&(0,0,1,0,0,0)&(1,1,0,1,0,0)\\
-432&(1,1,1,1,0)&(0,1,1,0,1)\\
13824&(1,0,0,0)&(0,1,1,0)\\
27&(1,0,1,0)&(0,1,0,1)\\
-4608&(1,0,1,0,0)&(1,1,0,1,0).
\end{array}
\tag{A.186.16}
\]

\emph{Equality with the Selmer image and the rank conclusion.}
The preceding arguments prove
\(\operatorname{Sel}^2(J_c/\Q)\subseteq\ker M_c\)
under the explicitly specified global coordinates.
The rational two-torsion has dimension one, since
\(xF_c(x)\) has exactly two irreducible factors;
this is also the branch-divisor calculation in (A.186.6),
or \cite[Lemma 4.3(3), p.251]{LegStollDescent2001}.
A.182 proves for each curve an actual rational divisor
class of infinite order and that all rational torsion
has order at most two.  Hence the rank is at least one,
the torsion is exactly \(\Z/2\), and finite generation
makes \(\dim J_c(\Q)/2J_c(\Q)\geq2\).
The finite Kummer exact sequence
\[
0\longrightarrow J_c(\Q)/2J_c(\Q)
 \longrightarrow\operatorname{Sel}^{2}(J_c/\Q)
 \longrightarrow\Sha(J_c/\Q)[2]\longrightarrow0
\tag{A.186.17}
\]
now forces equality of all three two-dimensional groups
\(J_c(\Q)/2J_c(\Q)\), \(\operatorname{Sel}^{2}\),
and \(\ker M_c\).  It proves rank one and
\(\Sha(J_c/\Q)[2]=0\).  Any nonzero element of
two-power order would have a nonzero multiple of order
two, so in fact \(\Sha(J_c/\Q)[2^\infty]=0\).
Applying the finite Kummer exact sequence at \(2^n\)
proves (A.186.2), with its actual transition maps.
Finally the isomorphisms (A.186.4) and the isogenies
already proved in A.182--A.185 give
\(\operatorname{rank}J(C_8)(\Q)=8+6=14\).

\emph{The actual generators of each completed Mordell--Weil group.}
For each of the following rational points \(P_c\), put
\(D_c=[P_c-\infty]\) and
\(T_c=[(0,0)-\infty]\).  A binary word in the table
means the coordinate vector, in the listed order of
\(G_c\) in (A.186.13), of the quartic squareclass
in (A.186.7):
\[
\begin{array}{r|c|c|c}
c&P_c&e(D_c)&e(T_c)\\ \hline
-54&(3,9)&\mathtt{01110}&\mathtt{11101}\\
648&(6,108)&\mathtt{110100}&\mathtt{111100}\\
-432&(6,72)&\mathtt{11110}&\mathtt{01101}\\
13824&(24,2880)&\mathtt{0110}&\mathtt{1110}\\
27&(3,18)&\mathtt{1111}&\mathtt{0101}\\
-4608&(-8,64)&\mathtt{10100}&\mathtt{11010}.
\end{array}
\tag{A.186.18}
\]
These are precisely the monic-coordinate images of the
rational points in (A.182.7), under (A.186.4).
For complete squareclass witnesses, write
\(P_c=(a_c,b_c)\).  The following elements satisfy
\[
 a_c-\theta=R_{D,c}^{\,2}\prod_i g_i^{e(D_c)_i},
 \qquad
 -\theta=R_{T,c}^{\,2}\prod_i g_i^{e(T_c)_i}
 \quad\text{in }K_c^\times:
\tag{A.186.19}
\]
\[
\begin{array}{r|c|c}
c&R_{D,c}&R_{T,c}\\ \hline
-54&-\theta^3/9-\theta^2/3+\theta+3
    &-\theta^3/9-\theta^2/3+\theta+3\\
648&-\theta^3/108+\theta^2/18&-\theta^3/108\\
-432&-\theta^3/72+\theta^2/6-\theta/2
     &-\theta^3/72+\theta/2\\
13824&-\theta^3/288+\theta/2+6&-\theta^3/576\\
27&-\theta^3/18+\theta/2&-\theta^3/9\\
-4608&-\theta^2/24+2&-\theta^3/16+\theta^2/2-4\theta+34.
\end{array}
\tag{A.186.20}
\]
Direct multiplication modulo \(\theta^4+c\) verifies
all twelve identities, including the square factor
at the prime five in the \(c=13824\) row.  Both
columns of every row in (A.186.18) lie in the kernel
already proved in (A.186.16), and they are independent.
Thus \(D_c,T_c\) generate \(J_c(\Q)/2J_c(\Q)\).
The rank-one free quotient of \(J_c(\Q)\) is
infinite cyclic; its coordinate of \(D_c\) is odd
by that independence.  Multiplication by this odd
integer is invertible on \(\Z_2\), while \(T_c\)
is the full torsion subgroup.  The literal map
\[
 \Z_2\oplus\Z/2\longrightarrow\widehat{J_c(\Q)}_2,
          \qquad(a,e)\longmapsto aD_c+eT_c
\tag{A.186.21}
\]
is therefore a continuous bijective homomorphism.
Its source and target are compact Hausdorff groups,
so its inverse is continuous as well.  Reducing
(A.186.21) modulo \(2^n\) and applying (A.186.2)
also specifies the exact finite Selmer generators
at every level.

The complete local computation is
\path{code/point3_fermat8_genus2_rank_certificate.gp},
executed with local PARI/GP 2.19.1.  Its retained output
contains all twenty-nine literal generators, their
full certified coordinate matrices, the seventeen
point witnesses, all eighty-five products of local
Hilbert symbols, and the six exact kernels.  It makes
no external submission.  The earlier exploratory
search for these rows is retained separately; the
final certificate checks the displayed data directly
and invokes full flag-zero class/unit certification.
The generator coordinates and all twelve literal
square-root witnesses are retained separately in
\path{code/point3_fermat8_genus2_generator_coordinates.gp}
and its executed output.
\end{proof}
